% TOP_PROBLEM: {"problemId": "problem.beilinson-soule-vanishing-conjecture-for-rational-motivic-cohomology", "canonicalTitle": "Beilinson-Soule Vanishing Conjecture for Rational Motivic Cohomology", "releaseRank": 186, "primaryDomain": "algebra_representation_category", "primaryDomainLabel": "Algebra, representation & category theory", "displayStatus": "open_with_solved_subcases", "releaseStatus": "open_with_solved_subcases"}
% !TeX program = lualatex
% Definition and sources only; this file is not an LLM solution attempt.
\documentclass[11pt]{article}
\usepackage[margin=25mm]{geometry}
\usepackage{fontspec}
\setmainfont{Latin Modern Roman}
\usepackage{amsmath,amssymb,mathtools}
\usepackage{unicode-math}
\setmathfont{Latin Modern Math}
\usepackage{xurl,hyperref}
\hypersetup{colorlinks=true,urlcolor=blue}
\setcounter{secnumdepth}{0}
\setlength{\parindent}{0pt}
\setlength{\parskip}{5pt}
\setlength{\emergencystretch}{3em}
\begin{document}
\section*{\texorpdfstring{Beilinson-Soule Vanishing Conjecture for Rational Motivic Cohomology}{Beilinson-Soule Vanishing Conjecture for Rational Motivic Cohomology}}\label{186}
\subsection{Definitions and mathematical statement}
For a smooth scheme $X$ over a field, the selected rational motivic groups are
\[
 H^{n,i}(X)=\operatorname{gr}_{\gamma}^{i}K_{2i-n}(X)\otimes{\mathbb{Q}}.
\]
Here $K_j$ is algebraic $K$-theory and the gamma filtration is the filtration defined by the lambda-ring operations, with its rational graded pieces providing the weight decomposition. The Beilinson--Soul\'e assertion is
$H^{n,i}(X)=0$ whenever $n\le0$ and $(n,i)\ne(0,0)$, with the conventional zero groups outside the available $K$-theoretic weights. The exceptional group $H^{0,0}$ is not required to vanish. This is a rational statement about all smooth schemes over fields, not the stronger claim that the corresponding integral groups have no torsion.
\par\smallskip\textit{Formulation and concept references:} \href{https://arxiv.org/pdf/2408.06233}{[S1]}.

\subsection{Short English statement}
Do all nonpositive-degree rational motivic cohomology groups vanish except for the ordinary degree-zero, weight-zero part?
\subsection{Sources}
\begin{itemize}
\item[S1]F. Deglise, "Generic motives and motivic cohomology of fields", arXiv:2408.06233 [math.AG], v1 dated 12 August 2024.
\url{https://arxiv.org/pdf/2408.06233}.
\textit{Formulation source.}
\end{itemize}
\end{document}
